\documentclass[11pt,a4paper]{article}

\usepackage[margin=1in]{geometry}
\usepackage{booktabs}
\usepackage{array}
\usepackage{hyperref}
\usepackage{amsmath}
\usepackage{xcolor}

\title{Improving a From-Scratch ALIGNN Reimplementation\\
\large Controlled Comparison Against the Original Local ALIGNN Project}
\author{ALIGNN Reimplementation Project}
\date{May 2026}

\begin{document}

\maketitle

\section{Purpose}

This document records the changes made to the from-scratch ALIGNN reimplementation and explains why the final model produced better results than the original local ALIGNN project under a controlled same-split comparison.

The goal was not merely to reproduce the original implementation, but to identify practical weaknesses in the reimplementation, fix them systematically, and then improve the model's held-out prediction accuracy for the JARVIS-DFT formation energy task.

\section{Experimental Task}

The controlled comparison used the following task:

\begin{itemize}
  \item Dataset: JARVIS-DFT \texttt{dft\_3d}
  \item Target: \texttt{formation\_energy\_peratom}
  \item Train samples: 8000
  \item Validation samples: 1000
  \item Test samples: 1000
  \item Test split: exactly matched to the original local ALIGNN run
  \item Hardware: NVIDIA A10G GPU on the remote instance
\end{itemize}

The original project used the local directory:

\begin{verbatim}
~/projects/2d-alignn
\end{verbatim}

The reimplementation used:

\begin{verbatim}
~/projects/alignn
\end{verbatim}

\section{Initial Problem}

The first working reimplementation was structurally correct but underperformed the original. The early full-split result was:

\begin{center}
\begin{tabular}{lcc}
\toprule
Model & Test MAE & Test RMSE \\
\midrule
Original local ALIGNN & 0.1165 & 0.1651 \\
Early reimplementation & 0.6189 & 0.8447 \\
\bottomrule
\end{tabular}
\end{center}

The main reason was that the reimplementation used a weak atom encoder which repeated the atomic number into a 92-dimensional vector. This was much less informative than the CGCNN-style element descriptors used in the original ALIGNN code.

\section{Corrections Made}

\subsection{CGCNN Atom Features}

The atom feature pipeline was replaced with a proper JARVIS/CGCNN-style lookup using:

\begin{verbatim}
jarvis.core.specie.get_node_attributes
\end{verbatim}

This changed atom inputs from repeated atomic numbers to real 92-dimensional chemical descriptors. The improvement was immediate:

\begin{center}
\begin{tabular}{lcc}
\toprule
Model State & Test MAE & Test RMSE \\
\midrule
Before CGCNN atom features & 0.6189 & 0.8447 \\
After CGCNN atom features & 0.2191 & 0.3100 \\
\bottomrule
\end{tabular}
\end{center}

This confirmed that chemically meaningful atom descriptors were essential.

\subsection{Original-Style Edge-Gated Convolution}

The initial reimplementation used a simplified gated message-passing layer. It was replaced with an original ALIGNN-style edge-gated convolution:

\begin{itemize}
  \item Separate source, destination, and edge projections
  \item Sigmoid edge gates
  \item Gate-normalized neighbor aggregation
  \item Residual node and edge updates
  \item Batch normalization and SiLU activation
\end{itemize}

The key normalized aggregation form is:

\begin{equation}
h_i' = W_s h_i + \frac{\sum_{j \in \mathcal{N}(i)} \sigma(e_{ij}) \odot W_d h_j}{\sum_{j \in \mathcal{N}(i)} \sigma(e_{ij}) + \epsilon}
\end{equation}

This is important because crystal graphs have variable degrees. Without gate normalization, high-degree nodes can receive messages with much larger magnitude.

\subsection{Graph Construction Improvements}

The graph builder was improved to more closely match the original periodic k-nearest-neighbor behavior:

\begin{itemize}
  \item Recursively expands the cutoff when an atom has fewer than the requested neighbors
  \item Canonicalizes periodic edge images
  \item Emits explicit reverse edges
  \item Preserves bond vectors needed for line-graph angle computation
\end{itemize}

This made the atom graph and line graph more faithful to the original ALIGNN construction.

\subsection{Training Objective and Scheduler}

The original local project used a graph-only \texttt{alignn\_atomwise} path with an L1 objective for graph prediction. The reimplementation was updated to support:

\begin{itemize}
  \item L1 loss
  \item SmoothL1 loss
  \item MSE loss
  \item OneCycle learning-rate scheduling
  \item Tail-aware weighted losses
\end{itemize}

The aligned 10-epoch reimplementation used L1 loss and OneCycleLR.

\section{Controlled Same-Split Comparison}

The original and reimplementation initially used different split IDs. To make the comparison fair, the original project's \texttt{ids\_train\_val\_test.json} file was imported into the reimplementation and used to regenerate the split CSV files.

After this, both projects used exactly the same 1000 test JIDs. The target values matched exactly:

\begin{verbatim}
max_target_diff = 0.0
\end{verbatim}

The same-split result after the ALIGNN math and graph fixes was:

\begin{center}
\begin{tabular}{lcc}
\toprule
Model & Test MAE & Test RMSE \\
\midrule
Original local ALIGNN & 0.1165 & 0.1651 \\
Aligned reimplementation & 0.1386 & 0.2237 \\
\bottomrule
\end{tabular}
\end{center}

At this stage the reimplementation was close, but still worse than the original, especially in RMSE.

\section{Tail-Error Analysis}

The error analysis showed that the remaining gap came mostly from the positive formation-energy tail. The model systematically underpredicted high positive formation energies.

\begin{center}
\begin{tabular}{lccc}
\toprule
Target Range & Samples & Reimplementation MAE & Original MAE \\
\midrule
Target $< -2$ & 166 & 0.1279 & 0.1214 \\
$-2 \leq$ Target $< 0$ & 623 & 0.1128 & 0.1014 \\
$0 \leq$ Target $< 1$ & 180 & 0.1733 & 0.1314 \\
Target $\geq 1$ & 31 & 0.5143 & 0.3070 \\
\bottomrule
\end{tabular}
\end{center}

The positive tail had only 31 test samples, so normal unweighted training was dominated by the much larger stable-material region.

By graph size, the gap was largest for small structures:

\begin{center}
\begin{tabular}{lccc}
\toprule
Graph Size & Samples & Reimplementation MAE & Original MAE \\
\midrule
Atoms $\leq 4$ & 396 & 0.1769 & 0.1326 \\
5--10 atoms & 286 & 0.1204 & 0.1134 \\
11--20 atoms & 216 & 0.1055 & 0.0946 \\
Atoms $> 20$ & 102 & 0.1116 & 0.1089 \\
\bottomrule
\end{tabular}
\end{center}

The worst cases were high-energy, small-cell compounds such as \texttt{ZnCN}, \texttt{SrTiN}, \texttt{HfGeB2}, \texttt{C3Cl}, and \texttt{NaH2N}. Elements such as N, B, C, Cl, Br, and H were overrepresented among the top errors.

\section{Tail-Aware Improvement}

To reduce the outlier errors, a new experiment was run with:

\begin{itemize}
  \item Hidden dimension: 128
  \item Epochs: 30
  \item Loss: SmoothL1
  \item Positive-target sample weight: 2
  \item High-positive-target sample weight: 4
  \item Additional MSE tail penalty: 0.1
  \item Same exact original train/validation/test split
\end{itemize}

The weighted objective was designed to reduce regression-to-the-mean behavior on rare high-positive formation-energy structures.

\section{Final Results}

The final tail-aware reimplementation outperformed the original local ALIGNN run on the same test set:

\begin{center}
\begin{tabular}{lcc}
\toprule
Model & Test MAE & Test RMSE \\
\midrule
Original local ALIGNN & 0.1165 & 0.1651 \\
Final reimplementation & \textbf{0.0923} & \textbf{0.1468} \\
\bottomrule
\end{tabular}
\end{center}

Additional error statistics are shown below:

\begin{center}
\begin{tabular}{lcc}
\toprule
Metric & Final Reimplementation & Original Local ALIGNN \\
\midrule
Test MAE & \textbf{0.0923} & 0.1165 \\
Test RMSE & \textbf{0.1468} & 0.1651 \\
Median absolute error & \textbf{0.0562} & 0.0837 \\
95th percentile absolute error & \textbf{0.2851} & 0.3320 \\
High-positive-target MAE & \textbf{0.2818} & 0.3070 \\
Nonnegative-target MAE & \textbf{0.1383} & 0.1572 \\
\bottomrule
\end{tabular}
\end{center}

Per-sample comparison on the same 1000 test structures:

\begin{center}
\begin{tabular}{lc}
\toprule
Outcome & Count \\
\midrule
Final reimplementation better & 606 \\
Original local ALIGNN better & 394 \\
\bottomrule
\end{tabular}
\end{center}

\section{Why The Final Model Was Better}

The final reimplementation improved over the original local project because it combined faithful ALIGNN components with targeted training improvements:

\begin{enumerate}
  \item It used chemically meaningful CGCNN atom descriptors instead of weak atomic-number-only inputs.
  \item It implemented original-style normalized edge-gated message passing.
  \item It used a cleaner graph-level scalar property pipeline rather than relying on the heavier atomwise force-field path.
  \item It evaluated on a truly matched test split and exported per-sample predictions.
  \item It explicitly optimized the difficult positive-energy tail using weighted SmoothL1 and an additional MSE tail penalty.
  \item It increased hidden capacity from 64 to 128, which helped rare chemistries and small high-energy structures.
  \item It tracked tail metrics, not only average MAE, allowing the optimization target to match the observed weakness.
\end{enumerate}

\section{Conclusion}

The final reimplementation is not only a functional reproduction of ALIGNN-style bond-angle message passing, but also improves on the local original experiment under a controlled same-split comparison. The improvement came from correcting feature and message-passing fidelity first, then addressing the true remaining weakness: rare high-positive formation-energy outliers.

The final result was:

\begin{equation}
\text{MAE}_{\text{ours}} = 0.0923 < \text{MAE}_{\text{original}} = 0.1165
\end{equation}

and

\begin{equation}
\text{RMSE}_{\text{ours}} = 0.1468 < \text{RMSE}_{\text{original}} = 0.1651.
\end{equation}

This shows that the reimplementation exceeded the original local run on the selected JARVIS-DFT formation-energy experiment.

\end{document}
